\documentclass[11pt]{article}
\usepackage[utf8]{inputenc}
\usepackage[T1]{fontenc}
\usepackage[spanish]{babel}
\usepackage{amsmath, amssymb, amsthm, mathtools}
\usepackage{braket}
\usepackage{dsfont}
\usepackage{graphicx}
\usepackage{float}
\usepackage[margin=2.5cm]{geometry}
\usepackage{hyperref}
\hypersetup{hidelinks}
\emergencystretch=2em
\allowdisplaybreaks

\newcommand{\ii}{\mathrm{i}}
\newcommand{\ee}{\mathrm{e}}
\newcommand{\Id}{\mathds{1}}
\newcommand{\Om}{\Omega_{99}^{\mathrm{ef}}}

\title{Un fon\'on en un bloque polarit\'onico}
\author{Ricardo Del Rosal Gardu\~no \\ \small Doctorado en F\'isica, BUAP}
\date{Notas de doctorado}

\begin{document}
\maketitle

\tableofcontents
\bigskip

\section{Idea de la reducci\'on}

En la nota de din\'amica del modelo HC se obtuvo la evoluci\'on general del
hamiltoniano efectivo en el espacio desplazado. Para cada bloque \'optico
$n>0$ y cada n\'umero efectivo polarit\'on-fon\'on $m>0$, el hamiltoniano
conecta \'unicamente los dos estados
\begin{equation}
  \ket{+,n,m-1},
  \qquad
  \ket{-,n,m}.
\end{equation}
Esto ocurre porque el t\'ermino de interacci\'on efectivo
\begin{equation}
  -\frac{\hbar\nu\lambda}{2}
  \left(
  \hat b\hat S^{(n)\dagger}
  +
  \hat b^\dagger\hat S^{(n)}
  \right)
\end{equation}
destruye un fon\'on y sube el polarit\'on, o crea un fon\'on y baja el
polarit\'on. En ambos procesos se conserva
\begin{equation}
  \boxed{
  \hat M^{(n)}
  =
  \hat b^\dagger\hat b+\hat S^{(n)\dagger}\hat S^{(n)}.
  }
\end{equation}

La diferencia con la nota anterior es que ahora no se considera una
distribuci\'on coherente de muchos fonones. Se toma un \'unico estado de
n\'umero fon\'onico y, por tanto, el sistema queda dentro de un solo doblete.
El manuscrito elige
\begin{equation}
  \boxed{
  n=9,
  \qquad
  m=9.
  }
\end{equation}
As\'i, el subespacio din\'amico relevante es
\begin{equation}
  \boxed{
  \mathcal H_{99}
  =
  \operatorname{span}
  \left\{
  \ket{+,9,8},
  \ket{-,9,9}
  \right\}.
  }
\end{equation}
La primera etiqueta despu\'es del signo polarit\'onico es el bloque \'optico
$n=9$. La \'ultima etiqueta es el n\'umero de fonones. El estado
$\ket{+,9,8}$ tiene el polarit\'on superior y ocho fonones; el estado
$\ket{-,9,9}$ tiene el polarit\'on inferior y nueve fonones. Ambos tienen
$M^{(9)}=9$.

\section{Estado inicial dentro del doblete}

La forma general de un doblete $(n,m)$ es
\begin{equation}
  \ket{\psi_{nm}(0)}
  =
  a_{nm}\ket{+,n,m-1}
  +
  b_{nm}\ket{-,n,m},
\end{equation}
con
\begin{equation}
  \boxed{
  a_{nm}=\tilde c_{m-1}d_n^+,
  \qquad
  b_{nm}=\tilde c_md_n^-.
  }
\end{equation}
Los coeficientes $\tilde c_m$ vienen del estado vibracional desplazado y los
coeficientes $d_n^\pm$ vienen de la descomposici\'on polarit\'onica del estado
\'optico-electr\'onico inicial.

Para el bloque elegido,
\begin{equation}
  \boxed{
  \ket{\psi_{99}(0)}
  =
  a_{99}\ket{+,9,8}
  +
  b_{99}\ket{-,9,9}.
  }
\end{equation}
El manuscrito toma como estado inicial preparado el estado inferior del
doblete:
\begin{equation}
  \boxed{
  \ket{\psi_{99}(0)}=\ket{-,9,9}.
  }
\end{equation}
Esto equivale a fijar
\begin{equation}
  \boxed{
  a_{99}=0,
  \qquad
  b_{99}=1.
  }
\end{equation}
F\'isicamente, el sistema comienza con el polarit\'on inferior y con nueve
fonones. La \'unica transici\'on permitida por el hamiltoniano efectivo es
hacia $\ket{+,9,8}$, porque subir el polarit\'on consume un fon\'on.

\section{Par\'ametros del doblete}

Para un doblete general, la matriz efectiva tiene la forma
\begin{equation}
  \tilde H_{\mathrm{ef}}^{(n,m)}
  =
  \hbar A_{nm}\Id_{2\times2}
  +
  \frac{\hbar}{2}
  \begin{pmatrix}
  \delta_n & -k_m\\
  -k_m & -\delta_n
  \end{pmatrix},
\end{equation}
donde
\begin{equation}
  A_{nm}
  =
  \Omega n+\frac{\nu\lambda^2}{4}
  +
  \nu\left(m-\frac12\right),
\end{equation}
\begin{equation}
  \delta_n=2g\sqrt n-\nu,
  \qquad
  k_m=\nu\lambda\sqrt m.
\end{equation}
La frecuencia de Rabi efectiva del doblete es
\begin{equation}
  \Omega_{nm}^{\mathrm{ef}}
  =
  \sqrt{\delta_n^2+k_m^2}.
\end{equation}

Al sustituir $n=9$ y $m=9$ se obtiene
\begin{equation}
  \boxed{
  A_{99}
  =
  9\Omega+\frac{\nu\lambda^2}{4}
  +
  \frac{17}{2}\nu.
  }
\end{equation}
La desintonizaci\'on polarit\'on-fon\'on queda
\begin{equation}
  \boxed{
  \delta_9
  =
  2g\sqrt9-\nu
  =
  6g-\nu.
  }
\end{equation}
El acoplamiento del doblete es
\begin{equation}
  \boxed{
  k_9
  =
  \nu\lambda\sqrt9
  =
  3\nu\lambda.
  }
\end{equation}
En las gr\'aficas se usa con frecuencia el par\'ametro
\begin{equation}
  \boxed{
  \eta=\nu\lambda,
  \qquad
  k_9=3\eta.
  }
\end{equation}
Con esta notaci\'on, la frecuencia efectiva es
\begin{equation}
  \boxed{
  \Omega_{99}^{\mathrm{ef}}
  =
  \sqrt{(6g-\nu)^2+9\nu^2\lambda^2}
  =
  \sqrt{(6g-\nu)^2+9\eta^2}.
  }
\end{equation}

La condici\'on de resonancia del bloque $n=9$ es
\begin{equation}
  \boxed{
  \delta_9=0
  \qquad\Longleftrightarrow\qquad
  \nu=6g.
  }
\end{equation}
Si $\nu$ se aleja de $6g$, la transferencia entre las dos ramas ya no es
completa: la oscilaci\'on conserva la misma frecuencia efectiva, pero su
amplitud se reduce por el factor $k_9/\Omega_{99}^{\mathrm{ef}}$.

\section{Evoluci\'on temporal}

La soluci\'on general en el doblete es
\begin{equation}
  \ket{\psi_{99}(t)}
  =
  p_{99}^{+}(t)\ket{+,9,8}
  +
  p_{99}^{-}(t)\ket{-,9,9}.
\end{equation}
De la soluci\'on exacta de la nota de din\'amica,
\begin{align}
  p_{nm}^{+}(t)
  &=
  \ee^{-\ii A_{nm}t}
  \left[
  \left(
  \cos\frac{\Omega_{nm}^{\mathrm{ef}}t}{2}
  -
  \ii\frac{\delta_n}{\Omega_{nm}^{\mathrm{ef}}}
  \sin\frac{\Omega_{nm}^{\mathrm{ef}}t}{2}
  \right)a_{nm}
  \right.
  \notag\\
  &\hspace{3.2cm}
  \left.
  +
  \ii\frac{k_m}{\Omega_{nm}^{\mathrm{ef}}}
  \sin\frac{\Omega_{nm}^{\mathrm{ef}}t}{2}
  b_{nm}
  \right],
  \\
  p_{nm}^{-}(t)
  &=
  \ee^{-\ii A_{nm}t}
  \left[
  \ii\frac{k_m}{\Omega_{nm}^{\mathrm{ef}}}
  \sin\frac{\Omega_{nm}^{\mathrm{ef}}t}{2}
  a_{nm}
  \right.
  \notag\\
  &\hspace{3.2cm}
  \left.
  +
  \left(
  \cos\frac{\Omega_{nm}^{\mathrm{ef}}t}{2}
  +
  \ii\frac{\delta_n}{\Omega_{nm}^{\mathrm{ef}}}
  \sin\frac{\Omega_{nm}^{\mathrm{ef}}t}{2}
  \right)b_{nm}
  \right].
\end{align}
Como en este caso $a_{99}=0$ y $b_{99}=1$, queda
\begin{equation}
  \boxed{
  p_{99}^{+}(t)
  =
  \ee^{-\ii A_{99}t}
  \ii\frac{k_9}{\Omega_{99}^{\mathrm{ef}}}
  \sin\!\left(
  \frac{\Omega_{99}^{\mathrm{ef}}t}{2}
  \right).
  }
\end{equation}
La amplitud de permanecer en la rama inferior es
\begin{equation}
  \boxed{
  p_{99}^{-}(t)
  =
  \ee^{-\ii A_{99}t}
  \left[
  \cos\!\left(
  \frac{\Omega_{99}^{\mathrm{ef}}t}{2}
  \right)
  +
  \ii\frac{\delta_9}{\Omega_{99}^{\mathrm{ef}}}
  \sin\!\left(
  \frac{\Omega_{99}^{\mathrm{ef}}t}{2}
  \right)
  \right].
  }
\end{equation}
El factor $\ee^{-\ii A_{99}t}$ es una fase com\'un. No afecta las
probabilidades, pero s\'i forma parte del estado vectorial completo.

\section{Probabilidades de las ramas polarit\'onicas}

La probabilidad de encontrar el sistema en la rama superior es
\begin{align}
  P_+(t)
  &=
  |p_{99}^{+}(t)|^2
  \notag\\
  &=
  \left|
  \ee^{-\ii A_{99}t}
  \ii\frac{k_9}{\Omega_{99}^{\mathrm{ef}}}
  \sin\!\left(
  \frac{\Omega_{99}^{\mathrm{ef}}t}{2}
  \right)
  \right|^2.
\end{align}
La fase global y el factor $\ii$ tienen m\'odulo uno. Por tanto,
\begin{equation}
  \boxed{
  P_+(t)
  =
  \left(
  \frac{k_9}{\Omega_{99}^{\mathrm{ef}}}
  \right)^2
  \sin^2\!\left(
  \frac{\Omega_{99}^{\mathrm{ef}}t}{2}
  \right).
  }
\end{equation}

Para la rama inferior,
\begin{align}
  P_-(t)
  &=
  |p_{99}^{-}(t)|^2
  \notag\\
  &=
  \left|
  \cos\!\left(\frac{\Omega_{99}^{\mathrm{ef}}t}{2}\right)
  +
  \ii
  \frac{\delta_9}{\Omega_{99}^{\mathrm{ef}}}
  \sin\!\left(\frac{\Omega_{99}^{\mathrm{ef}}t}{2}\right)
  \right|^2.
\end{align}
El primer t\'ermino es real y el segundo es puramente imaginario, de modo que
no aparece t\'ermino cruzado:
\begin{equation}
  P_-(t)
  =
  \cos^2\!\left(\frac{\Omega_{99}^{\mathrm{ef}}t}{2}\right)
  +
  \left(
  \frac{\delta_9}{\Omega_{99}^{\mathrm{ef}}}
  \right)^2
  \sin^2\!\left(\frac{\Omega_{99}^{\mathrm{ef}}t}{2}\right).
\end{equation}
Como
\begin{equation}
  (\Omega_{99}^{\mathrm{ef}})^2=\delta_9^2+k_9^2,
\end{equation}
se tiene
\begin{equation}
  \frac{\delta_9^2}{(\Omega_{99}^{\mathrm{ef}})^2}
  =
  1-
  \frac{k_9^2}{(\Omega_{99}^{\mathrm{ef}})^2}.
\end{equation}
Sustituyendo,
\begin{equation}
  \boxed{
  P_-(t)
  =
  1-
  \left(
  \frac{k_9}{\Omega_{99}^{\mathrm{ef}}}
  \right)^2
  \sin^2\!\left(
  \frac{\Omega_{99}^{\mathrm{ef}}t}{2}
  \right).
  }
\end{equation}
Esto confirma que $P_+(t)+P_-(t)=1$ dentro del doblete. No hay p\'erdida de
probabilidad porque se est\'a resolviendo un subespacio cerrado del
hamiltoniano efectivo.

\section{Inversi\'on polarit\'onica}

La inversi\'on polarit\'onica mide si el sistema se encuentra preferentemente
en la rama superior o en la inferior. Para el bloque $n=9$ se define
\begin{equation}
  \boxed{
  \mathcal W(t)
  =
  \bra{\psi_{99}(t)}
  \hat S_z^{(9)}
  \ket{\psi_{99}(t)}.
  }
\end{equation}
El operador $\hat S_z^{(9)}$ act\'ua como
\begin{equation}
  \hat S_z^{(9)}\ket{+,9}=+\ket{+,9},
  \qquad
  \hat S_z^{(9)}\ket{-,9}=-\ket{-,9}.
\end{equation}
Los estados vibracionales $\ket8$ y $\ket9$ son ortogonales, de modo que los
t\'erminos cruzados desaparecen. Entonces
\begin{align}
  \mathcal W(t)
  &=
  |p_{99}^{+}(t)|^2
  -
  |p_{99}^{-}(t)|^2
  \notag\\
  &=
  P_+(t)-P_-(t).
\end{align}
Usando las probabilidades anteriores,
\begin{equation}
  \boxed{
  \mathcal W(t)
  =
  -1+
  2
  \left(
  \frac{k_9}{\Omega_{99}^{\mathrm{ef}}}
  \right)^2
  \sin^2\!\left(
  \frac{\Omega_{99}^{\mathrm{ef}}t}{2}
  \right).
  }
\end{equation}
Al inicio, $\mathcal W(0)=-1$, porque el sistema fue preparado en
$\ket{-,9,9}$. Si la transferencia es completa, la inversi\'on puede alcanzar
$+1$, lo que corresponde al estado $\ket{+,9,8}$.

\subsection{Caso resonante}

En resonancia $\delta_9=0$, por lo que
\begin{equation}
  \Omega_{99}^{\mathrm{ef}}=|k_9|.
\end{equation}
La inversi\'on se simplifica a
\begin{equation}
  \mathcal W_R(t)
  =
  -1+2\sin^2\!\left(\frac{k_9t}{2}\right).
\end{equation}
Usando la identidad
\begin{equation}
  1-2\sin^2\left(\frac{x}{2}\right)=\cos x,
\end{equation}
se obtiene
\begin{equation}
  \boxed{
  \mathcal W_R(t)
  =
  -\cos(k_9t).
  }
\end{equation}
Esta expresi\'on muestra una oscilaci\'on completa entre $-1$ y $+1$. La
frecuencia est\'a fijada por $k_9=3\eta$. Si $\eta=0.01g$, entonces
$k_9=0.03g$.

\section{Inversi\'on electr\'onica}

La inversi\'on electr\'onica se define como
\begin{equation}
  W(t)
  =
  \bra{\psi_{99}(t)}
  \hat\sigma_z
  \ket{\psi_{99}(t)}.
\end{equation}
Aunque el sistema oscila entre dos ramas polarit\'onicas, cada polarit\'on
contiene la misma cantidad de poblaci\'on electr\'onica excitada y base. En el
bloque $n=9$,
\begin{equation}
  \ket{+,9}
  =
  \frac{\ket{e,8}+\ket{g,9}}{\sqrt2},
  \qquad
  \ket{-,9}
  =
  \frac{\ket{e,8}-\ket{g,9}}{\sqrt2}.
\end{equation}
Para cada una de estas ramas,
\begin{equation}
  \bra{+,9}\hat\sigma_z\ket{+,9}=0,
  \qquad
  \bra{-,9}\hat\sigma_z\ket{-,9}=0.
\end{equation}
En principio, $\hat\sigma_z$ tiene elementos cruzados entre $\ket{+,9}$ y
$\ket{-,9}$. Sin embargo, en el estado del doblete esos polaritones est\'an
acompa\~nados por estados fon\'onicos distintos:
\begin{equation}
  \ket{+,9,8},
  \qquad
  \ket{-,9,9}.
\end{equation}
Como $\braket{8|9}=0$, los t\'erminos cruzados no contribuyen. Por tanto,
\begin{equation}
  \boxed{
  W(t)=0.
  }
\end{equation}
La lectura f\'isica es importante: hay intercambio polarit\'on-fon\'on, pero no
aparece inversi\'on electr\'onica neta. El movimiento ocurre entre ramas
vestidas que, vistas electr\'onicamente, tienen poblaciones balanceadas.

\clearpage
\section{Gr\'aficas}

En esta secci\'on se incluyen las gr\'aficas asociadas a la din\'amica de un
fon\'on en un bloque polarit\'onico. En las comparaciones se contrasta la
evoluci\'on generada por el hamiltoniano HC completo con la soluci\'on
anal\'itica del hamiltoniano efectivo. Se usa $\eta=0.01g$ y se toma el bloque
$n=9$ como referencia resonante, para el cual $\nu=6g$.

\begin{figure}[H]
  \centering
  \includegraphics[width=0.95\linewidth]{\detokenize{../../notas-doctorado/imgs/Un fonon en bloque polaritonico/InversionPolaritonica.png}}
  \caption{Inversi\'on polarit\'onica calculada a partir del resultado
  anal\'itico del hamiltoniano efectivo con $\eta=0.01g$ y $\nu=5.99g$. Al estar
  muy cerca de la resonancia del bloque $n=9$, la oscilaci\'on conserva una
  amplitud grande, aunque la desintonizaci\'on evita que el caso sea
  exactamente resonante.}
\end{figure}

\begin{figure}[H]
  \centering
  \includegraphics[width=0.95\linewidth]{\detokenize{../../notas-doctorado/imgs/Un fonon en bloque polaritonico/Comparacion-IP-R.png}}
  \caption{Comparaci\'on de la inversi\'on polarit\'onica entre la soluci\'on
  num\'erica del hamiltoniano HC completo y la soluci\'on anal\'itica del
  hamiltoniano efectivo en el caso resonante. Para $\eta=0.01g$, las curvas se
  superponen dentro de la escala de la figura, lo que indica que el modelo
  efectivo captura correctamente la oscilaci\'on polarit\'on-fon\'on.}
\end{figure}

\begin{figure}[H]
  \centering
  \includegraphics[width=0.95\linewidth]{\detokenize{../../notas-doctorado/imgs/Un fonon en bloque polaritonico/Comparacion-IE.png}}
  \caption{Comparaci\'on de la inversi\'on electr\'onica. El resultado
  anal\'itico del hamiltoniano efectivo predice $W(t)=0$. La soluci\'on
  num\'erica del hamiltoniano completo reproduce ese comportamiento tanto en
  resonancia como fuera de resonancia, porque las ramas polarit\'onicas tienen
  poblaci\'on electr\'onica balanceada.}
\end{figure}

\begin{figure}[H]
  \centering
  \includegraphics[width=0.95\linewidth]{\detokenize{../../notas-doctorado/imgs/Un fonon en bloque polaritonico/Comparacion-IP-NR1.png}}
  \caption{Comparaci\'on de la inversi\'on polarit\'onica en un caso no
  resonante con $\eta=0.01g$ y $\nu=5.9g$. Para el bloque $n=9$, esta frecuencia
  produce $\delta_9=6g-\nu=0.1g$, de modo que la amplitud de transferencia se
  reduce por el factor $k_9/\Omega_{99}^{\mathrm{ef}}$.}
\end{figure}

\begin{figure}[H]
  \centering
  \includegraphics[width=0.95\linewidth]{\detokenize{../../notas-doctorado/imgs/Un fonon en bloque polaritonico/Comparacion-IP-NR2.png}}
  \caption{Comparaci\'on de la inversi\'on polarit\'onica para un caso no
  resonante construido de otra manera. Se mantiene $\nu=6g$, que resonar\'ia
  con el bloque $n=9$, pero se mide el bloque $n=10$. Como la resonancia de
  ese bloque requerir\'ia $\nu=2g\sqrt{10}$, la din\'amica observada queda
  desintonizada.}
\end{figure}

\begin{figure}[H]
  \centering
  \includegraphics[width=0.95\linewidth]{\detokenize{../../notas-doctorado/imgs/Un fonon en bloque polaritonico/Fidelidad-R.png}}
  \caption{Fidelidad entre los estados evolucionados con el hamiltoniano HC
  completo y con el hamiltoniano efectivo en el caso resonante
  $(\eta=0.01g,\nu=6g)$. La fidelidad permanece muy cercana a uno, lo que
  muestra que los t\'erminos descartados en la reducci\'on efectiva tienen un
  efecto peque\~no en esta escala de tiempo.}
\end{figure}

\begin{figure}[H]
  \centering
  \includegraphics[width=0.95\linewidth]{\detokenize{../../notas-doctorado/imgs/Un fonon en bloque polaritonico/Fidelidad-NR1.png}}
  \caption{Fidelidad para el caso no resonante con $\nu=5.9g$. La comparaci\'on
  cuantifica qu\'e tan cerca permanece la evoluci\'on efectiva de la evoluci\'on
  completa cuando el doblete $n=9,m=9$ ya no est\'a en resonancia exacta.}
\end{figure}

\begin{figure}[H]
  \centering
  \includegraphics[width=0.95\linewidth]{\detokenize{../../notas-doctorado/imgs/Un fonon en bloque polaritonico/Fidelidad-NR2.png}}
  \caption{Fidelidad para el caso no resonante donde se mantiene $\nu=6g$ pero
  se observa el bloque $n=10$. Este caso separa la idea de resonancia global de
  la resonancia por bloque: la frecuencia fon\'onica puede estar ajustada a un
  bloque y, al mismo tiempo, estar fuera de resonancia para otro.}
\end{figure}

\section{Lectura f\'isica}

El caso de un solo fon\'on en un bloque polarit\'onico deja ver la estructura
m\'as simple de la din\'amica efectiva. No hay suma sobre muchos valores de
$m$, no hay colapso por de-fasaje entre frecuencias distintas y no aparece una
distribuci\'on fon\'onica extendida. Todo se reduce a una oscilaci\'on entre
dos estados:
\begin{equation}
  \ket{-,9,9}
  \longleftrightarrow
  \ket{+,9,8}.
\end{equation}
En resonancia, la transferencia puede ser completa. Fuera de resonancia, la
frecuencia efectiva aumenta, pero la amplitud de transferencia disminuye.
La inversi\'on polarit\'onica detecta directamente esa oscilaci\'on. La
inversi\'on electr\'onica, en cambio, permanece nula porque cada polarit\'on
tiene igual peso de estado electr\'onico excitado y base, y los t\'erminos de
coherencia electr\'onica quedan acompa\~nados por estados fon\'onicos
ortogonales.

\end{document}
